% ============================================================
%  Section 9 — Limitations and next experiments
% ============================================================
\section{Limitations and Next Experiments}
\label{sec:roadmap}

The project is a course prototype with useful evidence and important gaps. We
prioritize the work by the evidence needed to make the company-adaptation claim
reproducible.

\begin{enumerate}[leftmargin=*, itemsep=3pt]
  \item \textbf{Broaden SkillOpt evaluation.} The real library now runs through
        the product API on balanced Company B TRAIN and SELECTION cohorts. The
        next experiment should use the full held-out set, repeat with multiple
        seeds, preserve all candidate traces, and evaluate a selected skill
        once on TEST before claiming a company-wide gain.
  \item \textbf{Confirm the targeted coding result.} Evaluate v003 on the full
        Company B SELECTION and held-out TEST sets, with model calls paced to
        stay within provider limits. Include code accuracy, citation support,
        JSON parse rate, provider errors, and paired uncertainty.
  \item \textbf{Strengthen retrieval evaluation.} The current rule corpus has
        only a few chunks per company, making Recall@5 saturated. Re-chunk at
        rule level, compare fixed and heading-based strategies on the same
        queries and vector model, and report per-query results.
  \item \textbf{Improve task quality.} Expand the synthetic corpus, audit
        company-rule and gold-label consistency, and add independently labeled
        real scans before making real-world accuracy claims. Keep arithmetic
        validation deterministic and report error-type precision and recall.
  \item \textbf{Harden deployment.} Add authentication, role-based review,
        secret management, backups, retention controls, and monitoring before
        using real company financial documents.
\end{enumerate}

No post-baseline score projections are reported. The next release should be
judged by repeatable held-out measurements rather than a target inferred from
unrelated benchmarks.
